\section{The native launch at a registered finite receiver}\label{sec:receiver}
Hold the three source profiles in \eqref{eq:launch} at their constant source-off endpoints, with p zero. The following result concerns that prepared plateau. Active later RX or WT evolution would require its actual histories and is not silently frozen by the mathematics.

\subsection{Placement and the measured component}
Choose a same-width receiving copy with the already defined proper frame
\begin{equation}\label{eq:receiverplacement}
Q=B(\pi/2)=[\ey,\ez,\ex],\qquad
O_R=O+(R-wK_c)\ez.
\end{equation}
Its local circle 0 has
\begin{equation}\label{eq:nativecontour}
X=O_R+wK_cQ\ey=O+R\ez,\qquad
q(\theta)=X+b(\cos\theta\,\ez+\sin\theta\,\ex),\quad
b=a=wR_\Gamma,\quad u=\ey.
\end{equation}
This rigid test placement uses a published orientation but not a published positive-radius H/R ring position. It changes neither the shell shape nor the loop radius. Circle-center X and body-center \(O_R\) differ, as Figure~\ref{fig:native} shows.

\begin{figure}[tbp]
\centering\resizebox{\textwidth}{!}{\input{figures/native}}
\caption{Geometry of the native-launch test. Left: the three source circles have distinct centers; their horizontal projections are radial segments because their planes are vertical. The arrows are the fixed oriented normals, not signs of the launched means. Right: the actual receiving circle lies in the xz plane with increasing theta clockwise in this view, normal \(+\ey\), and \(b=wR_\Gamma\). The long R separation is broken schematically; the picture is not to scale. The body center is displaced by \(wK_c\) from the circle center.}\label{fig:native}
\end{figure}

Let \(\ell_0=wK_c\) and \(J_*=\sum_i|j_i|\). If \(R>\ell_0+a+b\), the whole receiving disk is separated from every filled source disk by at least
\begin{equation}\label{eq:clearance}
\delta_D\ge R-\ell_0-a-b>0.
\end{equation}
The full copies and their circles are also separated by enclosing balls in the range \(R\ge100w\). All three distinct source centers remain in the exact C6 integrals.

Their aggregate leading moment about O is \(\pi a^2\alpha W_0\). At \(X=O+R\ez\), the leading curl field has y-component
\begin{equation}\label{eq:leadingcomponent}
h_{\rm lead}=-\frac{g_oa^2\alpha W_{0y}}{4R^3},\qquad
W_{0y}=\frac{57}{400000}>0.
\end{equation}
The centered \(O(r^{-5})\) error of one loop cannot be reused unchanged about O: translations of the three source centers can give an \(O(R^{-4})\) correction.

\begin{theorem}[A signed finite-radius native readout]\label{thm:native}
For \(\alpha\ne0\), placement \eqref{eq:receiverplacement}--\eqref{eq:nativecontour}, and every \(R\ge100w\),
\begin{equation}\label{eq:thirdbound}
|S_b-h_{\rm lead}|<\tfrac13|h_{\rm lead}|.
\end{equation}
For positive alpha the launch reading is negative; its conjugate is exactly the opposite and positive. The radius remains \(b=wR_\Gamma\).
\end{theorem}
\begin{proof}
The estimate has three contributions:
\begin{align}
|S_b-h_{\rm lead}|&\le\mathcal E(R),\label{eq:Ebound}\\
\mathcal E(R)&=g_oJ_*\left[
\frac{6a^2\ell_0}{(R-\ell_0)^4}
+\frac{3a^4}{2(R-\ell_0-a)^5}
+\frac{3a^2b^2}{2(R-\ell_0-a-b)^5}
\right].\notag
\end{align}
For the translation term, direct differentiation of \eqref{eq:kernelG} gives
\begin{equation}\label{eq:DG}
D_vG(x)=|x|^{-4}
\{-15(n\cdot v)nn^T+3(n\cdot v)I+3(nv^T+vn^T)\},
\end{equation}
so \(\norm{D_vG(x)}_{\rm op}\le24|v|/|x|^4\). The segment from \(X-O\) to \(X-Y_i\) stays outside radius \(R-\ell_0\); multiply by \(g_o|j_i|a^2/4\) and sum. The centered remainder \eqref{eq:remainder} gives the second term. Disk averaging \eqref{eq:diskaveragebound}, with \eqref{eq:clearance}, gives the third. Equation~\eqref{eq:Ebound} bounds the \emph{area-averaged} reading relative to the central leading component; it is not a pointwise error bound replacing all disk points by the center.

From the exact launched weights,
\begin{equation}\label{eq:Jratio}
\frac{J_*}{|\alpha|W_{0y}}=\frac{57002\sqrt3}{85500}<\frac65,\qquad
0<K_c<1,\quad \tfrac12<R_\Gamma<1.
\end{equation}
Put \(q=w/R\le1/100\). Dividing \eqref{eq:Ebound} by \(|h_{\rm lead}|\) yields
\begin{equation}\label{eq:rationalbound}
\frac{\mathcal E(R)}{|h_{\rm lead}|}<
\frac{144q}{5(1-q)^4}
+\frac{36q^2}{5}\{(1-2q)^{-5}+(1-3q)^{-5}\}.
\end{equation}
Each term increases on this interval. At \(q=1/100\) the right side is
\begin{equation}\label{eq:rationalvalue}
\frac{7804605427237675733225000}{25890315358991719010197977}
=0.301448835945798\ldots<\frac13.
\end{equation}
This proves \eqref{eq:thirdbound}. The exact integral is linear in all three means; conjugation reverses them all and hence reverses the full finite-radius reading.
\end{proof}
The range is sufficient and configuration-specific, not an optimal threshold. This native placement does not generally have zero V at X; the zero-center result belongs to the separate coaxial control.

\begin{table}[tbp]\centering\small
\begin{tabular}{@{}rrrr@{}}\toprule
\(R/w\)&Finite contour \(S_b\)&\(h_{\rm lead}\)&\(\mathcal E(R)\)\\\midrule
100&\(-1.0251976124210\,10^{-11}\)&\(-1.0252192872890\,10^{-11}\)&\(2.2207690784546\,10^{-12}\)\\
200&\(-1.2815173359258\,10^{-12}\)&\(-1.2815241091113\,10^{-12}\)&\(1.3656205101897\,10^{-13}\)\\\bottomrule
\end{tabular}
\caption{Saved numerical corroboration at \(g_o=\alpha=w=1\), using all three source circles. Angular resolutions 64, 128 and 256 and a disk integration of H agree in the GM06 evidence. The all-distance guarantee is the analytic bound \eqref{eq:rationalbound}--\eqref{eq:rationalvalue}.}\label{tab:native}
\end{table}

\subsection{The 1,024-sample statement has its own premises}
Suppose N additional equally spaced points on the receiving rim have exact signed vector records \(v_k=B_k^TV(q(\theta_k))\). Reassemble the vectors before combining frames:
\begin{equation}\label{eq:SN}
S_{b,N}=\frac{2}{bN}\sum_{k=0}^{N-1}t(\theta_k)\cdot B_kv_k,\qquad
\theta_k=2\pi k/N.
\end{equation}
If \(F(\theta)=V(q(\theta))\cdot t(\theta)\) is \(C^2\), the composite trapezoidal error is
\begin{equation}\label{eq:quadratureerror}
|S_{b,N}-S_b|\le\frac{2\pi^2}{3bN^2}\sup|F''|.
\end{equation}
Using \(q'=bt,q''=-b(e_1\cos\theta+e_2\sin\theta)\), and the kernel derivatives,
\begin{equation}\label{eq:Fsecond}
\sup|F''|\le\frac{g_oaJ_*}{2}
\left(\frac{2b^2}{\delta^3}+\frac{3b}{\delta^2}+\frac1\delta\right)
\end{equation}
for contour-wire clearance delta. Indeed \(|F''|\le b^2M_{2,V}+3bM_{1,V}+M_{0,V}\), while inverse-distance derivatives are bounded by \(2/\delta^3,1/\delta^2,1/\delta\).

At \(R=100w\), use \(\delta>97w\), \(a=b>w/2\), \(a=b<w\), \eqref{eq:Jratio} and \(\pi^2<10\). For N=1024,
\begin{equation}\label{eq:1024}
\frac{|S_{b,1024}-S_b|}{|h_{\rm lead}|}
<
\frac{64000000}{1024^2}
\left(\frac2{97^3}+\frac3{97^2}+\frac1{97}\right)<\frac23.
\end{equation}
The continuous reading has magnitude greater than \(2|h_{\rm lead}|/3\), so the discrete sum retains its sign. These are new contour samples, not the old array's body-center sites. The result is for \(R=100w\), not fixed N at every larger R. Tangent measurement error at most epsilon per sample adds at most \(2\epsilon/b\); no instrument precision is presumed. The one-sample coaxial corollary has a separate symmetry premise.

\subsection{Arrival and completion of the steady reading}\label{sec:delay}
For a source history that vanishes before \(t_0\) and is constant after the jump, define on the rim C
\begin{equation}\label{eq:delays}
\tau_{\min}^{C}=\min_{P\in C,i,s}\frac{|P-\gamma_i(s)|}{c_e},
\qquad
\tau_{\max}^{C}=\max_{P\in C,i,s}\frac{|P-\gamma_i(s)|}{c_e}.
\end{equation}
The contour is silent before \(t_0+\tau_{\min}^C\). Its whole static record is available after the strict time \(t>t_0+\tau_{\max}^C\), provided the sources remain constant. Intermediate records need their actual delayed histories; the earliest permitted contribution does not guarantee a nonzero aggregate at the first-arrival instant.

The reviewed timing correction is
\begin{equation}\label{eq:delaymaxequal}
\boxed{\tau_{\max}^{D_b}=\tau_{\max}^{C}.}
\end{equation}
For a fixed source point, every interior disk point is a convex combination of two rim points on a chord. Convexity of distance bounds its distance by the larger endpoint distance. Maximizing over source points proves equality. Disk and rim minima can differ. The strict post-maximum inequality gives a neighborhood where all delayed arguments lie in the smooth plateau, justifying Stokes and differentiation over the disk.

For coaxial circles the common latest delay is
\(\sqrt{z^2+(a+b)^2}/c_e\), while the rim and disk minima are respectively
\(\sqrt{z^2+(a-b)^2}/c_e\) and \(\sqrt{z^2+(a-b)_+^2}/c_e\).
For the native placement, sufficient bounds are
\begin{equation}\label{eq:delaybounds}
\tau_{\min}^{D_b}\ge\frac{R-\ell_0-a-b}{c_e},\qquad
\tau_{\max}^{D_b}=\tau_{\max}^{C}\le\frac{R+\ell_0+a+b}{c_e}.
\end{equation}
At \(R=100w\) the ceiling is \(101.82961419736\,w/c_e\). This is the adopted straight-delay clock, not a derived electromagnetic speed or native physical duration.
